%!TeX root=../tese.tex
%("dica" para o editor de texto: este arquivo é parte de um documento maior)
% para saber mais: https://tex.stackexchange.com/q/78101

\chapter{As packages \texttt{imegoodies} e \texttt{imelooks}}
\label{ann:imegoodlooks}

Este modelo inclui as \textit{packages} \texttt{imegoodies} e \texttt{imelooks},
que você pode querer usar em outros documentos \LaTeX.

\texttt{imegoodies} inclui um grande número de \textit{packages} que são
comumente usadas e bastante úteis. Em geral, você pode incluí-la em seus
documentos sem que isso cause problemas de compatibilidade. Se, no
entanto, algo não funcionar, você pode editar o arquivo para eliminar
a \textit{package} responsável pelo problema se ela não for necessária.
\texttt{imegoodies} ainda inclui vários comentários explicativos sobre as
\textit{packages} carregadas.

\texttt{imelooks} também inclui um grande número de \textit{packages}, mas
estas são relacionadas mais explicitamente à aparência do documento
(fontes, cores, margens etc.). Você também pode utilizá-la em outros
documentos se quiser se aproximar da aparência deste modelo. \texttt{imelooks}
reconhece diversos parâmetros que ativam/desativam aspectos específicos:

\begin{itemize}
  \item \texttt{fonts} carrega as fontes deste modelo (libertinus e
        sourcecodepro), além de outros pequenos ajustes relacionados.
        Esta opção é sempre ativada por padrão; para desativá-la, use
        \texttt{nofonts}

  \item \texttt{spacing} utiliza os espaçamentos definidos neste modelo (margens,
        espaço entre parágrafos, indentação da primeira linha do parágrafo
        etc.). Esta opção é sempre ativada por padrão; para desativá-la, use
        \texttt{nospacing}

  \item \texttt{captions} e \texttt{footnotes} fazem respectivamente as legendas
        (das figuras e tabelas) e as notas de rodapé de acordo com este modelo.
        Estas opções são sempre ativadas por padrão; para desativá-las, use
        \texttt{nocaptions} e \texttt{nofootnotes}

  \item \texttt{autohttp} acrescenta o prefixo \texttt{http://} a URLs criadas
        com \texttt{url} que não incluam o \textit{schema}. Esta opção é
        sempre ativada por padrão; para desativá-la, use \texttt{noautohttp}

  \item \texttt{hidelinks}, \texttt{borderlinks} e \texttt{colorlinks} definem a
        aparência dos hiperlinks. \texttt{hidelinks} faz os hiperlinks sem
        nenhuma formatação especial; \texttt{borderlinks} faz os hiperlinks
        serem envidos por um quadrado colorido (apenas na tela; o quadrado
        não é impresso); \texttt{colorlinks} faz o texto dos hiperlinks ser
        colorido. A opção \texttt{colorlinks} é sempre ativada por padrão

  \item \texttt{biblatex} carrega a \textit{package} \texttt{biblatex} e os
        estilos bibliográficos deste modelo. Esta opção é sempre ativada
        por padrão; para desativá-la, use \texttt{nobiblatex}
  \item \texttt{raggedbib} faz a bibliografia (com \texttt{biblatex}) ser
        formatada com alinhamento à esquerda ao invés de justificado.
        Esta opção é sempre ativada por padrão, exceto quando o estilo
        bibliográfico é \texttt{plainnat-ime} (usado nas teses); para
        desativá-la, use \texttt{noraggedbib}; para ativá-la incondicionalmente,
        use \texttt{raggedbib}
  \item \texttt{bibstyle=?} selectiona um estilo bibliográfico específico.
        O estilo padrão é \texttt{numeric}, exceto em pôsteres e apresentações
        (\texttt{beamer-ime}) e \textit{reports} (\texttt{plainnat-ime})

  \item \texttt{listings} carrega a \textit{package} \texttt{listings} e diversas
        configurações relacionadas usadas neste modelo. Esta opção é
        sempre ativada por padrão; para desativá-la, use \texttt{nolistings}

  \item \texttt{greeny}, \texttt{bluey}, \texttt{sandy} ativam esquemas de cores
        diferentes para pôsteres e apresentações (o padrão é \texttt{bluey})

  \item \texttt{beamer} \textbf{des}ativa algumas \textit{packages} que
        são incompatíveis com a classe \texttt{beamer} (note que as opções
        \texttt{slides} e \texttt{presentation}, discutidas abaixo, já fazem isso)

  \item \texttt{presentation} (ou \texttt{slides}) e \texttt{poster} ativam as
        opções relevantes para, respectivamente, apresentações com
        \texttt{beamer} ou pôsteres com \texttt{tcolorbox}

  \item \texttt{report} ativa as opções relevantes para documentos com
        capítulos (cabeçalhos das páginas, características do sumário etc.)

  \item \texttt{thesis} ativa a opção \texttt{report} e também define o que é
        necessário para a geração da capa das teses de acordo com este modelo

  \item \texttt{resumoabstract} define os comandos \texttt{resumo} e \texttt{abstract}
        de acordo com este modelo. Esta opção é ativada por padrão com
        \texttt{report}; para desativá-la, use \texttt{noresumoabstract}

  \item \texttt{brazilian} verifica se a língua portuguesa está ativa no
        documento e, em caso negativo, gera um erro. Esta opção é
        ativada por padrão com a opção \texttt{thesis}; para desativá-la,
        use \texttt{nobrazilian}
\end{itemize}
